\documentclass[10pt,a4paper]{amsart}
\usepackage{amsmath,amssymb,amsthm,graphicx}
% Retained preamble for the Stacks Project chapter snapshot in this directory.
% The chapter's own preamble.tex is a copy of the shared Stacks Project preamble; the block
% below is sliced MECHANICALLY, line for line, from it. Nothing was added, removed or
% redefined; the wrapper class, the xr-hyper/hyperref wiring and the externaldocument list
% are omitted because the assembler supplies the document class.
%
% lines taken from preamble.tex: 3, 5, 9, 11, 12, 156, 157, 158, ... (51 lines total)
%
% package coverage restored from the chapter's own preamble (2026-10-13): xy, xy-2cell

\usepackage[all]{xy}
\xyoption{2cell}
\UseAllTwocells

\newenvironment{reference}{\comment}{\endcomment}
\newenvironment{slogan}{\comment}{\endcomment}
\newenvironment{history}{\comment}{\endcomment}
\theoremstyle{plain}
\newtheorem{theorem}[subsection]{Theorem}
\newtheorem{proposition}[subsection]{Proposition}
\newtheorem{lemma}[subsection]{Lemma}
\theoremstyle{definition}
\newtheorem{definition}[subsection]{Definition}
\newtheorem{example}[subsection]{Example}
\newtheorem{exercise}[subsection]{Exercise}
\newtheorem{situation}[subsection]{Situation}
\theoremstyle{remark}
\newtheorem{remark}[subsection]{Remark}
\newtheorem{remarks}[subsection]{Remarks}
\numberwithin{equation}{subsection}
\def\lim{\mathop{\mathrm{lim}}\nolimits}
\def\colim{\mathop{\mathrm{colim}}\nolimits}
\def\Spec{\mathop{\mathrm{Spec}}}
\def\Hom{\mathop{\mathrm{Hom}}\nolimits}
\def\Ext{\mathop{\mathrm{Ext}}\nolimits}
\def\SheafHom{\mathop{\mathcal{H}\!\mathit{om}}\nolimits}
\def\SheafExt{\mathop{\mathcal{E}\!\mathit{xt}}\nolimits}
\def\Sch{\mathit{Sch}}
\def\Mor{\mathop{\mathrm{Mor}}\nolimits}
\def\Ob{\mathop{\mathrm{Ob}}\nolimits}
\def\Sh{\mathop{\mathit{Sh}}\nolimits}
\def\NL{\mathop{N\!L}\nolimits}
\def\CH{\mathop{\mathrm{CH}}\nolimits}
\def\proetale{{pro\text{-}\acute{e}tale}}
\def\etale{{\acute{e}tale}}
\def\QCoh{\mathit{QCoh}}
\def\Ker{\mathop{\mathrm{Ker}}}
\def\Im{\mathop{\mathrm{Im}}}
\def\Coker{\mathop{\mathrm{Coker}}}
\def\Coim{\mathop{\mathrm{Coim}}}
\DeclareMathSymbol{\boxtimes}{\mathbin}{AMSa}{"02}
\def\QCohstack{\mathcal{QC}\!\mathit{oh}}
\def\Cohstack{\mathcal{C}\!\mathit{oh}}
\def\Spacesstack{\mathcal{S}\!\mathit{paces}}
\def\Quotfunctor{\mathrm{Quot}}
\def\Hilbfunctor{\mathrm{Hilb}}
\def\Curvesstack{\mathcal{C}\!\mathit{urves}}
\def\Polarizedstack{\mathcal{P}\!\mathit{olarized}}
\def\Complexesstack{\mathcal{C}\!\mathit{omplexes}}
\def\Pic{\mathop{\mathrm{Pic}}\nolimits}
\def\Picardstack{\mathcal{P}\!\mathit{ic}}
\def\Picardfunctor{\mathrm{Pic}}
\def\Deformationcategory{\mathcal{D}\!\mathit{ef}}
\usepackage{cleveref}
\crefname{originallemma}{original lemma}{original lemmas}
\Crefname{originallemma}{Original lemma}{Original lemmas}
\makeatletter
\newlabel{algebra-lemma-ass-zero-divisors}{{00LD}{source}{}{sourcepointer-1}{}}
\newlabel{algebra-lemma-ass-zero-divisors@cref}{{[originallemma][0][]00LD}{[1][0][]source}{}{}{}}
\makeatother
\crefname{originallemma}{original lemma}{original lemmas}
\Crefname{originallemma}{Original lemma}{Original lemmas}
\makeatletter
\newlabel{algebra-lemma-bourbaki}{{0312}{source}{}{sourcepointer-2}{}}
\newlabel{algebra-lemma-bourbaki@cref}{{[originallemma][0][]0312}{[1][0][]source}{}{}{}}
\makeatother
\crefname{originallemma}{original lemma}{original lemmas}
\Crefname{originallemma}{Original lemma}{Original lemmas}
\makeatletter
\newlabel{algebra-lemma-flat-localization}{{00HT}{source}{}{sourcepointer-3}{}}
\newlabel{algebra-lemma-flat-localization@cref}{{[originallemma][0][]00HT}{[1][0][]source}{}{}{}}
\makeatother
\title{Stacks, Discriminants: lemma-explain-condition}
\author{The Stacks Project Authors}
\date{stacks-project:discriminant.tex:local-snapshot-20261012; GNU FDL 1.2 or later}
\begin{document}
\maketitle
\noindent\textit{Editorial note.} Statement span is lines 3187-3206 of sources/stacks-discriminant/pinned-original-source.tex; proof span is lines 3208-3243, the complete author proof (34 body lines).\par\medskip
\noindent\textit{Editorial note.} Navigation stubs: 3 cited-but-not-included labels are supplied as non-load-bearing pointers with their pinned tags where the snapshot tag index has them.\par\medskip
\noindent\textit{Editorial note.} extra\_packages is empty on purpose: the retained preamble already loads xy with options and a second declaration clashes in the assembler wrapper.\par\medskip
\noindent\textit{Editorial note.} Text is retained verbatim, including the author's own typos and placeholders.\par\medskip
\noindent\textit{Editorial note.} Queue index 10 of queue\_discriminant.json (pb=1731, stmt 3187-3206, proof 3208-3243).\par\medskip
\noindent\textit{Allowed standard prerequisites (editorial).}\begin{itemize}
\item Explicit prerequisite boundary. The cited inputs are prerequisites; the argument in the retained span is not replaced by any cited result.
\end{itemize}\medskip
\ifdefined\hypertarget\hypertarget{sourcepointer-1}{}\fi
\noindent\textit{Source reference.} Original lemma 00LD: sources/stacks-discriminant/pinned-original-source.tex, source lines 3243--3243. Cited by the retained statement or proof but not part of the collected spans; supplied only as a navigation pointer. The argument does not depend on its content.\par\medskip
\ifdefined\hypertarget\hypertarget{sourcepointer-2}{}\fi
\noindent\textit{Source reference.} Original lemma 0312: sources/stacks-discriminant/pinned-original-source.tex, source lines 3243--3243. Cited by the retained statement or proof but not part of the collected spans; supplied only as a navigation pointer. The argument does not depend on its content.\par\medskip
\ifdefined\hypertarget\hypertarget{sourcepointer-3}{}\fi
\noindent\textit{Source reference.} Original lemma 00HT: sources/stacks-discriminant/pinned-original-source.tex, source lines 3243--3243. Cited by the retained statement or proof but not part of the collected spans; supplied only as a navigation pointer. The argument does not depend on its content.\par\medskip
\begin{lemma}
\label{lemma-explain-condition}
Let $A \to B$ be a map of Noetherian rings. Consider the conditions
\begin{enumerate}
\item nonzerodivisors of $A$ map to nonzerodivisors of $B$,
\item (1) holds and $Q(A) \to Q(A) \otimes_A B$ is flat,
\item $A \to B_\mathfrak q$ is flat for every
$\mathfrak q \in \text{Ass}(B)$,
\item (3) holds and $A \to B_\mathfrak q$ is flat for every $\mathfrak q$
lying over an element in $\text{Ass}(A)$.
\end{enumerate}
Then we have the following implications
$$
\xymatrix{
(1) & (2) \ar@{=>}[l] \ar@{=>}[d] \\
(3) \ar@{=>}[u] & (4) \ar@{=>}[l]
}
$$
If going up holds for $A \to B$ then (2) and (4) are equivalent.
\end{lemma}

\begin{proof}
The horizontal implications in the diagram are trivial.
Let $S \subset A$ be the set of nonzerodivisors so that
$Q(A) = S^{-1}A$ and $Q(A) \otimes_A B = S^{-1}B$. Recall that
$S = A \setminus \bigcup_{\mathfrak p \in \text{Ass}(A)} \mathfrak p$
by Algebra, Lemma \ref{algebra-lemma-ass-zero-divisors}.
Let $\mathfrak q \subset B$ be a prime lying over $\mathfrak p \subset A$.

\medskip\noindent
Assume (2). If $\mathfrak q \in \text{Ass}(B)$ then
$\mathfrak q$ consists of zerodivisors, hence (1) implies
the same is true for $\mathfrak p$. Hence
$\mathfrak p$ corresponds to a prime of $S^{-1}A$.
Hence $A \to B_\mathfrak q$ is flat by our assumption (2).
If $\mathfrak q$ lies over an associated prime $\mathfrak p$
of $A$, then certainly $\mathfrak p \in \Spec(S^{-1}A)$ and the
same argument works.

\medskip\noindent
Assume (3). Let $f \in A$ be a nonzerodivisor. If $f$ were a zerodivisor
on $B$, then $f$ is contained in an associated prime $\mathfrak q$
of $B$. Since $A \to B_\mathfrak q$ is flat by assumption, we conclude that
$\mathfrak p$ is an associated prime of $A$ by
Algebra, Lemma \ref{algebra-lemma-bourbaki}. This would imply that
$f$ is a zerodivisor on $A$, a contradiction.

\medskip\noindent
Assume (4) and going up for $A \to B$. We already know (1) holds.
If $\mathfrak q$ corresponds to a prime of $S^{-1}B$ then $\mathfrak p$
is contained in an associated prime $\mathfrak p'$ of $A$. By going up
there exists a prime $\mathfrak q'$ containing $\mathfrak q$ and lying
over $\mathfrak p$. Then $A \to B_{\mathfrak q'}$ is flat by
(4). Hence $A \to B_{\mathfrak q}$ is flat as a localization.
Thus $A \to S^{-1}B$ is flat and so is $S^{-1}A \to S^{-1}B$, see
Algebra, Lemma \ref{algebra-lemma-flat-localization}.
\end{proof}

\end{document}
